% Outline for the sea-ice-model section of the VFS paper.

\section{Sea-ice thermodynamics and surface coupling}\label{sec:sea-ice}

\subsection{Model scope and state}
\begin{itemize}
  \item State clearly that the present model is a single-category, thermodynamic sea-ice component with no horizontal ice dynamics.
  \item Define the prognostic state: ice fraction, ice thickness, snow mass on sea ice, and sea-ice surface temperature.
  \item List key parameters: ice density, latent heat of fusion, formation thickness, melt temperature, bare-ice and snow albedos, roughnesses, and snow-cover/insulation scales.
  \item Explain ownership: sea ice owns its phase stores; ocean owns heat and salinity; the surface system only blends properties and routes explicit transfers.
\end{itemize}
\subsection{Boundary properties and snow cover}
\begin{itemize}
  \item Define snow-cover fraction as a bounded function of snow-water equivalent.
  \item Describe snow blending of bare-ice and snow albedo/roughness.
  \item Describe the snow insulation factor used to reduce ice growth under snow cover.
  \item Explain how ice fraction overlays the ocean fraction in the surface-fraction state.
\end{itemize}

\subsection{Freeze, melt, and salt exchange}
\begin{itemize}
  \item Define the local freezing condition using sea-surface temperature, salinity-dependent freezing temperature, and upward ocean heat loss.
  \item Explain conversion of available heat to ice mass using latent heat and the handling of fraction and thickness after growth or complete melt.
  \item Explain energy-limited melt: snow melts before ice, then meltwater returns freshwater and latent heat to the ocean.
  \item Explain brine rejection: salt excluded during freezing is returned to the ocean and diagnosed explicitly.
  \item List the transfer-ledger entries for snowfall, freeze, brine rejection, snowmelt, and ice melt.
\end{itemize}

\subsection{Coupling diagnostics}
\begin{itemize}
  \item Define the ocean-facing heat, freshwater, salt, and brine-rejection flux diagnostics.
  \item Explain the area/fraction weighting used in storage and flux calculations.
  \item State how these diagnostics enter the same-support surface conservation assessment.
\end{itemize}

\subsection{Verification and results to report}
\begin{table}[t]
\centering
\caption{Planned sea-ice verification cases.}
\label{tab:ice-verification}
\begin{tabular}{p{0.28\linewidth}p{0.35\linewidth}p{0.27\linewidth}}
\hline
Case & Required diagnostic & Planned paper result \\
\hline
Freezing event & Ice mass gain, ocean heat loss, brine salt return & Water/salt/energy closure table \\
Snowfall and melt sequence & Snow storage, snowmelt-first behavior, ocean freshwater flux & Time series and conservation residual \\
Snow-insulation sensitivity & Ice growth with bare versus snow-covered ice & Growth-rate comparison plot \\
Fractional-ice surface blend & Effective albedo, roughness, and ocean/ice fractions & Property-blending panel \\
Coupled surface gate & Maximum relative residual across all sea-ice transfers & Pass/fail summary with tolerance \\
\hline
\end{tabular}
\end{table}

\subsection{Known scope limits}
\begin{itemize}
  \item Do not claim a dynamic sea-ice model: advection, deformation, ridging, rheology, leads, melt ponds, thickness distributions, and ice momentum are absent.
  \item The current growth/melt rule is flux-limited; conductive multilayer ice/snow thermodynamics remains future work.
  \item Present results as phase-change and coupling-conservation verification, not a validation of observed sea-ice extent or thickness climatology.
\end{itemize}
